% Core theory, expanded from the repository foundation and audited 2026-09-19.
\section{Statistical experiment and intervention boundary}\label{statistical-experiment-and-intervention-boundary}

\subsection{Unit, timeline, and full history}\label{unit-timeline-and-full-history}

Initially suppose episodes are independent and identically distributed. An episode starts with task features and context \(H_1=X\), then generates

\[
O=(H_1,A_1,R_1,L_2,A_2,R_2,L_3,\ldots,A_T,R_T,L_{T+1}).
\]

All history and observation spaces are standard Borel, policies and kernels are measurable, and regular conditional laws exist. Action sets are finite. Whenever an optimizing action is selected, ties are resolved by a fixed measurable rule. Conditional regressions are defined on target-relevant support; arbitrary measurable versions may be assigned elsewhere because zero-occupancy contributions vanish.

Here \(H_t=(X,A_1,R_1,L_2,\ldots,A_{t-1},R_{t-1},L_t)\) is the \textbf{entire available history before routing at decision \(t\)}. Histories can contain text, tool responses, error indicators, elapsed time, cumulative cost, and previous model choices. No Markov assumption on a short state embedding is used. \(A_t\) is a finite-valued macro action: the choice of a frozen model/configuration for the next prespecified execution block. \(L_{t+1}\) contains the resulting generated output and environment observations, and \(R_t\) is the reward or utility accumulated in that block. The next history is the deterministic concatenation of these variables.

An execution block might be one model invocation, one tool-use cycle, or a prespecified subtask. Choose it before logging; changing the block definition changes the intervention. An intervention on routing does not automatically intervene on each generated token or each low-level tool action.

Let \(T<\infty\) be a fixed maximum number of decision opportunities. Early completion is an absorbing state: after the observed stopping time \(\tau\le T\), the only available action is a no-op with logging and target probability one, the state remains absorbed, and all subsequent rewards are zero. Emit a terminal success/failure reward at \(\tau\), or at \(T\) if the task reaches the cap. Thus

\[
G=\sum_{t=1}^{T}R_t
\]

can represent terminal task quality, cumulative resource use, or a prespecified utility. A cap-induced failure is part of the outcome definition. A crash that loses the final outcome is missing data and requires an additional observation model; it is not automatically absorption. An unbounded horizon is outside the theorems below.

\subsection{Macro actions are stochastic interventions on a fixed kernel}\label{macro-actions-are-stochastic-interventions-on-a-fixed-kernel}

Write the observed law as

\begin{equation}\label{eq:core1}
dP(O)=dP_1(H_1)\prod_{t=1}^{T}
b_t(A_t\mid H_t)\,dK_t(R_t,L_{t+1}\mid H_t,A_t),
\end{equation}

where \(b_t\) is the behavior/logging router. The kernel \(K_t\) integrates over the selected model's token randomness and subsequent tool/environment randomness. Schematically, for generated content \(Z_t\),

\begin{equation}\label{eq:core2}
K_t(dr,dl\mid h,a)
=\int E_t(dr,dl\mid h,a,z)\,M_{a,t}(dz\mid h).
\end{equation}

Model identity includes weight revision, quantization where applicable, tokenizer, prompt template, decoding settings, context handling, and memory serialization. The harness includes tools, permissions, stopping rules, reward measurement, and state restoration. If a model switch also changes these, define the action as that entire bundle or standardize them. A model-name string alone does not guarantee consistency.

A target regime \(\pi=(\pi_1,\ldots,\pi_T)\) replaces \(b_t\) by \(\pi_t\) while retaining \(P_1\) and \(K_t\). A deterministic DTR \(d\) is the special case \(\pi_t(a\mid h)=\mathbf{1}\{a=d_t(h)\}\). For now, \(\pi\) is fixed independently of the evaluation sample and does not vary along statistical submodels of \(P\).

\textbf{Macro-action implication.} If the intervention preserves \(K_t\), its likelihood ratio concerns model choice, not generated-token likelihoods. Token distributions of two models need not overlap with one another. When the logger sometimes chooses each model at each relevant history, its mixture already contains the trajectories generated by those models. This does not permit evaluating an unseen model, a new model revision, or a changed harness: each changes \(K_t\), rather than just \(b_t\).

\section{Identification assumptions and target value}\label{identification-assumptions-and-target-value}

The following assumptions identify a causal value, rather than merely an observed-law functional.

\begin{enumerate}
\def\labelenumi{\arabic{enumi}.}
\tightlist
\item
  \textbf{Well-defined intervention, consistency, and no interference.} Executing a routing history produces the potential observations for that same history and configuration. Separate episodes do not affect one another through shared writable memory, changing external systems, or shared resource contention unless those mechanisms are explicitly part of the design and estimand.
\item
  \textbf{Sequential exchangeability.} Conditional on the recorded pre-routing \(H_t\), the choice \(A_t\) is independent of the relevant future potential outcomes under supported continuations. Prospective routing randomization with correctly logged probabilities supplies this condition by design. Observational routing requires that all common causes of routing and future outcomes be captured. A human operator's unrecorded assessment or a hidden router feature can violate it.
\item
  \textbf{Sequential support.} For $P^{\pi}_{H_t}$-almost every history $h$, \(\pi_t(a\mid h)>0\) implies \(b_t(a\mid h)>0\). This is a population statement, not a claim that every long text history must repeat in a finite sample. Practical near-positivity remains a serious estimation problem.
\item
  \textbf{Stable task and transition law.} The evaluation target uses the same initial task distribution and conditional execution kernels as the logged experiment, or a separately justified transport model. Model updates, workload changes, judge drift, and a new task mix are not repaired merely by adjusting routing propensities.
\item
  \textbf{Outcome observation and integrability.} Rewards are observed as defined, and expectations exist. For later finite-sample results assume \(|R_t|\le r\), and for efficiency assume the influence function is square integrable.
\end{enumerate}

The target is

\begin{equation}\label{eq:core3}
v^\pi=E_{P^\pi}[G],\qquad
dP^\pi=dP_1\prod_{t=1}^{T}\pi_t\,dK_t.
\end{equation}

The g-formula originates in longitudinal causal inference, and DTRs describe the corresponding adaptive decision rules; the same value is the object of OPE. These are overlapping theoretical traditions, not competing estimands. \citep{robins1986new}, \citep{murphy2003optimal}, \citep{jiang2016dr}.

\begin{proposition}[G-computation and importance weighting]\label{core-result-1}
Define \(V_{T+1}^{\pi}=0\), and recursively

\[
Q_t^{\pi}(h,a)=E_P\{R_t+V_{t+1}^{\pi}(H_{t+1})\mid H_t=h,A_t=a\},
\] \begin{equation}\label{eq:core4}
V_t^{\pi}(h)=\sum_a\pi_t(a\mid h)Q_t^{\pi}(h,a).
\end{equation}

Then
\begin{equation}\label{eq:core5}
v^\pi=E_P\{V_1^{\pi}(H_1)\}
=E_P(W_T^{\pi}G)
=E_P\left\{\sum_{t=1}^{T}W_t^{\pi}R_t\right\},
\end{equation}

where

\begin{equation}\label{eq:core6}
W_0^{\pi}=1,\qquad
W_t^{\pi}=\prod_{s=1}^{t}
\frac{\pi_s(A_s\mid H_s)}{b_s(A_s\mid H_s)}.
\end{equation}
\end{proposition}
\begin{proof}
Consistency and exchangeability identify each intervened conditional kernel with \(K_t\); sequential support permits the change of measure. Iterated integration of \eqref{eq:core3} gives the backward recursion and its first equality in \eqref{eq:core5}. In the ratio of \eqref{eq:core3} to \eqref{eq:core1}, \(P_1\) and every \(K_t\) cancel, leaving \(W_T^{\pi}\). This proves the full-trajectory identity. For a reward measurable immediately after decision \(t\), integrating all later ratios gives conditional expectation one on target-relevant prefixes; zero-weight prefixes contribute zero. Thus \(E(W_T^{\pi}R_t)=E(W_t^{\pi}R_t)\). Summing proves \eqref{eq:core5}. Absorbed steps contribute ratio one and reward zero.
\end{proof}

For a deterministic rule, the numerator in \eqref{eq:core6} is the indicator of matching the rule. The product \(\prod_t1/b_t(A_t\mid H_t)\) by itself is \textbf{not} the weight for a specified policy; the target numerator is essential. Likewise, self-normalized or clipped importance sampling is a different estimator and does not inherit exact unbiasedness automatically.

The same factorization proves the macro-action implication in the macro-action construction above even if \(M_{a,t}\) and \(M_{a',t}\) have disjoint supports. It is sufficient that the behavior mixture chooses each needed macro action; one never divides one model's token probability by the other model's token probability.

\section{Fixed-policy influence function and estimation}\label{fixed-policy-influence-function-and-estimation}

Suppress superscript \(\pi\) in \(Q_t,V_t,W_t\). Define the uncentered score

\begin{equation}\label{eq:core7}
\phi^{\pi}(O;Q,b)
=V_1(H_1)+\sum_{t=1}^{T}W_t
\{R_t+V_{t+1}(H_{t+1})-Q_t(H_t,A_t)\}.
\end{equation}

\begin{theorem}[Efficient influence function in the full-history model]\label{core-result-2}
For a fixed target \(\pi\), in the unrestricted dominated model \eqref{eq:core1}, provided the expression is square integrable,

\begin{equation}\label{eq:core8}
D^{\pi}(O)=\phi^{\pi}(O;Q,b)-v^{\pi}
\end{equation}

is the efficient influence function (EIF). It remains efficient when the behavior probabilities are known and fixed by design, in the otherwise unrestricted model. This efficiency statement is relative to the full-history model. Additional true Markov restrictions can yield a smaller efficiency bound and different efficient estimators; using a state embedding does not establish such restrictions. \citep{kallus2020drl}.
\end{theorem}
\begin{proof}
A regular submodel has score decomposition
\[
S=S_1(H_1)+\sum_{t=1}^{T}S_{b,t}(A_t,H_t)
+\sum_{t=1}^{T}S_{K,t}(R_t,L_{t+1},H_t,A_t),
\]

with \(E S_1=0\), \(E(S_{b,t}\mid H_t)=0\), and \(E(S_{K,t}\mid H_t,A_t)=0\). Each residual in \eqref{eq:core7} has conditional mean zero given \(H_t,A_t\), so \(E D^{\pi}=0\).

Differentiating the g-formula along this submodel gives

\begin{equation}\label{eq:core9}
\dot v^{\pi}
=E[(V_1-v^{\pi})S_1]
+\sum_{t=1}^{T}E[W_t\{R_t+V_{t+1}\}S_{K,t}].
\end{equation}

There is no behavior-score derivative because \(\pi\) is fixed. In the term for a transition at \(t\), past rewards have zero covariance with its conditionally mean-zero score, and integrating future target transitions replaces future rewards by \(V_{t+1}\). Subtracting \(Q_t\) in \eqref{eq:core9} changes nothing because \(E(S_{K,t}\mid H_t,A_t)=0\).

The baseline component \(V_1-v^{\pi}\) and the weighted transition residuals form orthogonal martingale increments in the full-history filtration. Their covariances with all other score increments are zero. Hence \eqref{eq:core9} equals \(E(D^{\pi}S)\). These components lie in the corresponding baseline and transition tangent spaces, proving that \eqref{eq:core8} is the canonical gradient. Removing the behavior tangent spaces when \(b\) is known does not change this gradient, because it is already orthogonal to them.
\end{proof}

This is the established longitudinal/OPE influence-function construction, expressed for an agent-level treatment. It does not imply that a high-dimensional history regression will be accurate with a small benchmark.

\begin{theorem}[Exact drift and the scope of double robustness]\label{core-result-3}
Let \(\bar Q_t,\bar b_t\) be arbitrary fixed fitted functions with positive denominators wherever needed, let \(\bar V_t=\sum_a\pi_t\bar Q_t\), and let \(\bar W_t=\prod_{s\le t}\pi_s/\bar b_s\). Define \(\bar\phi\) by substituting these functions in \eqref{eq:core7}. Then the following exact identity holds whenever its terms are integrable:

\begin{equation}\label{eq:core10}
\begin{split}
E_P(\bar\phi)-v^{\pi}
=\sum_{t=1}^{T}E_P\Bigg[\bar W_{t-1}
\sum_a\pi_t(a\mid H_t)
\frac{\bar b_t(a\mid H_t)-b_t(a\mid H_t)}{\bar b_t(a\mid H_t)}
\{\bar Q_t(H_t,a)-Q_t^{\pi}(H_t,a)\}\Bigg].
\end{split}
\end{equation}
\end{theorem}
\begin{proof}
Put \(e_t=\bar Q_t-Q_t^{\pi}\) and \(d_t=\bar V_t-V_t^{\pi}=\sum_a\pi_te_t\), with \(d_{T+1}=0\). The true Bellman residual has conditional mean zero even when multiplied by \(\bar W_t\). Therefore

\[
E(\bar\phi)-v^{\pi}
=E(d_1)+\sum_t E\{\bar W_t(d_{t+1}-e_t(H_t,A_t))\}.
\]

Collect adjacent \(d_t\) terms to obtain

\[
\sum_t E\{\bar W_{t-1}d_t-\bar W_te_t(H_t,A_t)\}.
\]

Conditional on \(H_t\), the second term integrates over \(A_t\) under \(b_t\). Substituting \(d_t=\sum_a\pi_te_t\) gives \eqref{eq:core10}.
\end{proof}

Two familiar consequences are exact unbiasedness if every \(\bar b_t=b_t\), regardless of the fitted outcome functions, and exact unbiasedness if every \(\bar Q_t=Q_t^{\pi}\), regardless of the fitted behavior functions. Conditional unbiasedness under known randomization requires that the fitted functions be obtained independently of the scored episode, for example by sample splitting. These are algebraic population statements; estimation without splitting need not be unbiased in finite samples. \citep{bang2005dr}.

Equation \eqref{eq:core10} also vanishes for stagewise mixed patterns when, at each \(t\), either \(\bar b_t=b_t\) or \(\bar Q_t=Q_t^{\pi}\). \textbf{The latter is the true target-continuation regression, not just a correctly specified regression of an incorrectly constructed pseudo-outcome.} In ordinary backward fitting, errors in \(\bar V_{t+1}\) can propagate into \(\bar Q_t\). A naive recursive regression does not automatically converge to \(Q_t^{\pi}\) under every mixed pattern. We therefore call \eqref{eq:core7} a longitudinal DR/OPE score and do not claim that the ordinary backward-regression fitting algorithm has the stronger algorithmic sequential-double-robustness property. Such a claim requires a suitable fitting construction and proof. \citep{luedtke2017sdr}.

For a useful rate bound, define \(\mu_t(dh,a)=P(H_t\in dh)\pi_t(a\mid h)\). If \(\bar b_t\ge\epsilon>0\) on target-supported actions and \(\|\bar W_{t-1}\|_\infty\le C_{t-1}\), Cauchy--Schwarz gives

\begin{equation}\label{eq:core11}
|E(\bar\phi)-v^{\pi}|
\le\sum_{t=1}^{T}\frac{C_{t-1}}{\epsilon}
\|\bar b_t-b_t\|_{L_2(\mu_t)}
\|\bar Q_t-Q_t^{\pi}\|_{L_2(\mu_t)}.
\end{equation}

These are sufficient, deliberately strong overlap conditions. They are not needed in this form for identification. Large products of ratios can still make estimation impractical.

\subsection{A reproducible fitting algorithm}\label{a-reproducible-fitting-algorithm}

Split independent units into a fixed number of folds. For each held-out fold, fit all nuisance functions on other folds. For each fixed policy, work backward from \(\widehat V_{T+1}=0\), regress \(R_t+\widehat V_{t+1}(H_{t+1})\) on \((H_t,A_t)\), and set \(\widehat V_t=\sum_a\pi_t\widehat Q_t\). Use the \textbf{actual logged randomization probabilities} when available; do not replace them by an estimated classifier for convenience. Score each held-out episode with \eqref{eq:core7}, then average over all episodes. Model fitting, preprocessing, representation learning using outcomes, and tuning belong entirely within training data. Additional nested splitting within training can help control nuisance overfitting but is not a substitute for the convergence assumptions below.

For comparison, include the empirical mean in directly randomized target-policy arms, sequential g-computation, per-decision IPW, and the longitudinal DR score. Unsupported regression predictions are model extrapolations, not empirical identification.

\begin{theorem}[Cross-fitted asymptotic normality]\label{core-result-4}
Suppose \(T\) and the number of folds are fixed, every fold size divided by \(n\) tends to a strictly positive constant, evaluation episodes are i.i.d., \(\pi\) is fixed, \(0<\sigma_\pi^2=E[(D^{\pi})^2]<\infty\), and for every fold:

\begin{itemize}
\tightlist
\item
  the fitted score converges to \(\phi^{\pi}(O;Q,b)\) in \(L_2(P)\);
\item
  its population drift in \eqref{eq:core10} is \(o_p(n^{-1/2})\), for example by \eqref{eq:core11} and the corresponding product-rate conditions;
\item
  score integrability and variance consistency hold, for example through bounded rewards, fixed overlap constants, and suitably bounded fitted regressions.
\end{itemize}

Then
\begin{equation}\label{eq:core12}
\sqrt n(\widehat v^{\pi}-v^{\pi})
=n^{-1/2}\sum_{i=1}^{n}D^{\pi}(O_i)+o_p(1)
\ \rightsquigarrow\ N(0,\sigma_\pi^2).
\end{equation}
\end{theorem}
\begin{proof}
Within each fold, decompose the estimation error into the empirical mean of the true EIF, the empirical process of fitted-score minus true-score, and the fitted-score population drift. Conditional on training data, the empirical-process term has mean zero and variance bounded by its \(L_2(P)\) squared error divided by fold size. It is therefore \(o_p(n^{-1/2})\). A fixed finite number of folds preserves this order without requiring independence across fitted folds. The drift is negligible by assumption. The classical i.i.d. central limit theorem applies to the true EIF average.
\end{proof}

Estimate variance with the sample variance of held-out scores, and use \(\widehat v\pm1.96\widehat\sigma/\sqrt n\) for a pointwise approximate 95\% interval. With unknown propensities, both nuisance errors of order \(o_p(n^{-1/4})\) suffice for \eqref{eq:core11}, but high-dimensional agent histories do not guarantee that rate. With known propensities the drift is zero; efficiency still requires the score to converge to the EIF.

\textbf{Point consistency does not establish confidence-interval validity.} For example, suppose a behavior fit converges to a fixed incorrect \(\bar b\), while a fitted correct outcome model has error \(O_p(n^{-1/2})\). Equation \eqref{eq:core10} then generally has drift \(O_p(n^{-1/2})\), rather than \(o_p(n^{-1/2})\). The value estimate can be consistent, but the first-order randomness of fitting Q can enter its limiting distribution. A Wald interval computed solely from the evaluation-score variance omits that contribution. Cross-fitting removes evaluation-sample overfitting; it does not remove this first-order drift. Reporting all four nuisance-misspecification cells as though DR point consistency guaranteed 95\% coverage would be incorrect. To justify such inference, derive the additional nuisance influence terms, obtain a negligible drift by stronger conditions, or use a separately justified inferential method.

\textbf{Known-randomization corollary.} If \(b\) is known exactly and the cross-fitted scores converge in \(L_2(P)\) to a fixed \(\phi^{\pi}(O;\bar Q,b)\), then \eqref{eq:core10} gives zero drift even when \(\bar Q\ne Q^{\pi}\). The proof of Theorem~\ref{core-result-4} instead gives asymptotic normality with influence function \(\phi^{\pi}(O;\bar Q,b)-v^{\pi}\), provided it has a finite nonzero variance and the empirical variance is consistent. Score-based intervals are valid under these conditions, but efficiency is not guaranteed. This is a central reason to prefer prospective logging with known randomization probabilities in the first agent experiment.

\section{Task clusters, contrasts, and resource outcomes}\label{task-clusters-contrasts-and-resource-outcomes}

\subsection{Repeated seeds and branched trajectories}\label{repeated-seeds-and-branched-trajectories}

If each task is run with several seeds, branches, or models, the independent unit is usually the \textbf{task cluster}. Randomly assigning episodes from one task to both training and evaluation leaks task information and can understate uncertainty. Split by task and, where appropriate, by related task family or source repository.

For \(J\) independent and identically distributed task clusters with a fixed number \(m\) of valid root-to-terminal runs per task, assume each run has the specified behavior-law marginal and fit nuisances outside the entire held-out cluster. Define the task-weighted estimator as the average of cluster mean scores. Under the corresponding cluster-scale score-convergence, negligible-drift, and variance-consistency conditions, the cluster influence contribution is

\begin{equation}\label{eq:core13}
D_j^{\mathrm{cluster}}=m^{-1}\sum_{r=1}^{m}D^{\pi}(O_{jr}).
\end{equation}

Use the empirical variance of cluster scores divided by \(J\), not an episode-level variance divided by \(Jm\). Dependence from paired random seeds is allowed within a cluster. Selected continuation branches do not automatically have the root-episode marginal; clustering alone cannot repair that sampling bias. Their continuation estimand and sampling weights must be derived separately before a valid task-level contribution is averaged. Unequal or outcome-dependent numbers of runs require explicitly defining task versus episode weighting and rederiving the corresponding ratio or weighted estimator. A task-weighted average and an episode-weighted average need not estimate the same population.

For a fixed benchmark with repeated seeds, uncertainty over random execution seeds is conditional on that benchmark; it is not automatically uncertainty over a wider task population. If the behavior router learns from outcomes of previous episodes, the data may require martingale or adaptive-experiment theory. The i.i.d. theorem is not a justification for that setting.

\subsection{Policy contrasts and a cost vector}\label{policy-contrasts-and-a-cost-vector}

For two frozen policies evaluated on the same logs,

\begin{equation}\label{eq:core14}
\Delta=v^{\pi_1}-v^{\pi_0},\qquad
D_\Delta=D^{\pi_1}-D^{\pi_0}.
\end{equation}

Compute the variance of the paired score difference; summing two separate variances discards covariance. For \(J\) task clusters, apply the same operation to cluster means.

It is usually better to preserve a reward vector

\[
R_t^{\mathrm{vec}}=(R_t^{\mathrm{quality}},C_t^{\mathrm{tokens}},
C_t^{\mathrm{latency}},C_t^{\mathrm{money}},C_t^{\mathrm{switch}})
\]

and estimate its value vector and covariance. A prespecified scalar utility is a linear combination, such as quality minus \(\lambda^\top\)cost. State units and \(\lambda\); selecting \(\lambda\) after inspecting evaluation results introduces additional selection. Money can be near zero for local inference while latency and energy are not. Serial stage times sum to latency; parallel execution needs an explicit elapsed-time outcome rather than assuming stage-time additivity. Quantiles and tail-risk outcomes need their own influence-function arguments.

One can optimize quality subject to expected token and latency budgets using a finite candidate set and simultaneous uncertainty bounds. This is not the same as a per-task hard budget, which must be built into the policy's allowable actions and absorbing rule.

\section{Policy learning and a finite-class improvement guarantee}\label{policy-learning-and-a-finite-class-improvement-guarantee}

\subsection{Correct optimal continuation}\label{correct-optimal-continuation}

For an unrestricted supported policy class and the specified scalar utility, Bellman's recursion on full history is

\[
V_{T+1}^{*}=0,\qquad
Q_t^{*}(h,a)=E[R_t+V_{t+1}^{*}(H_{t+1})\mid H_t=h,A_t=a],
\] \begin{equation}\label{eq:core15}
V_t^{*}(h)=\max_{a\in\mathcal A_t(h)}Q_t^{*}(h,a),\qquad
d_t^{*}(h)\in\arg\max_aQ_t^{*}(h,a).
\end{equation}

The admissible set \(\mathcal A_t(h)\) must respect support and any hard resource restrictions. The proof is backward induction: the last decision maximizes its immediate expected remaining reward; conditional on the optimal continuation at later stages, maximizing \eqref{eq:core15} is optimal at the current stage. Full history suffices; a Markov state compression is optional and requires justification.

The naive \(\arg\max_aE[Y\mid H_t=h,A_t=a]\) usually evaluates \textbf{behavior-policy continuation}, not optimal continuation. For example, at decision 1, action 0 leads to terminal mean 0.6 regardless of decision 2. Action 1 leads to terminal reward 1 if decision 2 uses action 1 and 0 otherwise. If the behavior policy chooses that latter action only 10\% of the time, the observed conditional means are 0.6 versus 0.1, so the naive rule chooses 0. The optimal continuation after action 1 gives value 1, so the true optimal first choice is 1. Both decision-2 actions can have positive probability; this failure is not merely a positivity violation.

Practical policy training can use Q-learning, policy search, or existing routing algorithms. Independent evaluation is needed regardless of the training algorithm. Inference for the value of a frozen learned policy is easier than inference for the population optimum or the identity of an optimal rule near ties; no generic regularity claim about the latter is made here.

\begin{theorem}[Held-out finite-class improvement certificate]\label{core-result-5}
Condition on a training dataset that yields \(K\) candidate policies, a baseline \(\pi_0\), and fitted outcome regressions for all of them. Evaluate on \(n\) fresh i.i.d. episodes generated under \textbf{known} \(b\), and require sequential target support for every candidate and baseline. Suppose for every candidate and the baseline:

\[
|R_t|\le r,\quad
|\bar Q_t(h,a)|\le (T-t+1)r,\quad
\pi_t(a\mid h)/b_t(a\mid h)\le c_t
\]

on realized support, with deterministic finite constants. Set \(C_t=\prod_{s\le t}c_s\) and

\begin{equation}\label{eq:core16}
M=Tr+2r\sum_{t=1}^{T}C_t(T-t+1).
\end{equation}

Compute \(\bar\phi_k\) with known \(b\), and let

\begin{equation}\label{eq:core17}
\widehat\Delta_k=n^{-1}\sum_i(\bar\phi_k(O_i)-\bar\phi_0(O_i)),
\qquad
\epsilon_n=M\sqrt{\frac{8\log(2K/\alpha)}{n}}.
\end{equation}

With conditional probability at least \(1-\alpha\), simultaneously for all \(k\),

\begin{equation}\label{eq:core18}
|\widehat\Delta_k-(v^{\pi_k}-v^{\pi_0})|\le\epsilon_n.
\end{equation}

Consequently, choosing any candidate whose lower bound \(\widehat\Delta_k-\epsilon_n\) is positive, and retaining the baseline if none is positive, certifies positive expected-utility improvement whenever a switch is made, with error probability at most \(\alpha\).
\end{theorem}
\begin{proof}
The fitted score has conditional mean \(v^{\pi_k}\) by \eqref{eq:core10}, because \(b\) is known exactly. Its initial term is bounded by \(Tr\). The absolute residual at time \(t\) is at most \(r+(T-t)r+(T-t+1)r=2(T-t+1)r\), establishing \(|\bar\phi_k|\le M\). Each score difference lies in \([-2M,2M]\), an interval of width \(4M\). Hoeffding's inequality therefore gives

\[
P\{|\widehat\Delta_k-\Delta_k|>u\mid\mathrm{training}\}
\le2\exp\{-nu^2/(8M^2)\}.
\]

A union bound over \(K\) candidates and substitution of \eqref{eq:core17} prove \eqref{eq:core18}. Selection using simultaneous bounds preserves the event.
\end{proof}

The guarantee concerns expected utility under the specified task and harness distribution; it is not a safety guarantee for every task. It permits evaluation-set selection among a frozen finite class, but not inventing new candidates after looking at the evaluation set without accounting for that adaptation. With task clusters, use independent bounded cluster mean scores and replace \(n\) by the number of clusters. Multiple outcome/budget constraints require a union bound across both candidates and outcome channels. Estimated propensities introduce drift; the theorem then requires an explicit additional bias bound. Standard cross-fitting alone does not make this finite-sample statement exact.

This deliberately elementary bound is often conservative, especially with long horizons. Its purpose is to make the conditions for a valid claim transparent. Tighter empirical-Bernstein, confidence-sequence, or localized policy-class results are extensions, not established here. Safe policy improvement using DR OPE is already part of the prior literature, including \citep{jiang2016dr}.

\section{A support-aware routing intervention}\label{a-support-aware-routing-intervention}

Deterministic ``always use the large model'' rules can be nearly unsupported. An alternative is a bounded change in the odds of choosing that model where the logger already explores it. This is an application of established incremental propensity-score interventions, not a newly invented intervention family. \citep{kennedy2019incremental}.

For two models at active decision opportunities, write \(b_t(h)=P(A_t=1\mid H_t=h)\). Absorbed opportunities retain only their no-op, ratio one, and zero policy-derivative contribution. Let \(\delta_t(h)>0\) be a prespecified measurable routing-intensity function, for example \(\delta>1\) after a recoverable error and 1 elsewhere. Define

\begin{equation}\label{eq:core19}
q_t^{\delta}(1\mid h)
=\frac{\delta_t(h)b_t(h)}{1-b_t(h)+\delta_t(h)b_t(h)},
\qquad
Z_t(h)=1+\{\delta_t(h)-1\}b_t(h).
\end{equation}

This shifts odds, not probabilities by a fixed additive amount. It preserves deterministic routing where \(b=0\) or \(b=1\); therefore it cannot identify the benefit of trying an action never tried there. It identifies a different, supported question. If the intended action is literally ``switch versus continue,'' define a binary switch indicator relative to the previous model in \(H_t\), and apply the same construction to that action. Choosing model 1 and switching are not synonyms.

\begin{proposition}[Overlap control]\label{core-result-6}
For fixed \(\delta\), the realized likelihood ratios in \eqref{eq:core19} are

\begin{equation}\label{eq:core20}
\frac{q_t^{\delta}(A_t\mid H_t)}{b_t(A_t\mid H_t)}
=\frac{\delta_t(H_t)A_t+1-A_t}{Z_t(H_t)}.
\end{equation}

If \(0<\delta_{\min}\le\delta_t(h)\le\delta_{\max}<\infty\), each ratio is at most \(\max\{1,\delta_{\max},1/\delta_{\min}\}\). Structural zeros cause no realized infinite ratio. Products may still grow exponentially with \(T\).
\end{proposition}
\begin{proof}
\(Z_t\) is a convex combination of 1 and \(\delta_t\), hence lies between their minimum and maximum. Apply this to numerator \(\delta_t\) for \(A_t=1\) and 1 for \(A_t=0\).
\end{proof}

There are three different statistical targets:

\begin{enumerate}
\def\labelenumi{\arabic{enumi}.}
\tightlist
\item
  \textbf{Known randomized reference.} If \(b\) is known and fixed by design and \eqref{eq:core19} is defined from it, \(q^{\delta}\) is a fixed policy in that model. Theorems~\ref{core-result-2}--\ref{core-result-5} apply directly.
\item
  \textbf{A learned but frozen reference.} If a reference router \(b^{\mathrm{ref}}\) is learned on independent training data and frozen, the target is \(v^{q^{\delta}[b^{\mathrm{ref}}]}\), conditional on training. Use its fixed numerator and the actual evaluation logging denominator \(b\). The target need not equal \(v^{q^{\delta}[b]}\), and the ratio bound \eqref{eq:core20} is not automatic when the two references differ. Check support explicitly.
\item
  \textbf{The unknown behavior-adaptive population target.} If the target functional itself is \(\Psi_{\delta}(P)=v^{q^{\delta}[b_P]}(P)\), perturbing \(P\) perturbs the target policy. Treating that policy as fixed gives the wrong EIF in general.
\end{enumerate}

For case 3, in the unrestricted model, the corrected EIF is

\begin{equation}\label{eq:core21}
\begin{split}
D_{\delta}^{\mathrm{adaptive}}(O)
=D^{q^{\delta}}_{\mathrm{fixed}}(O)
+\sum_{t=1}^{T}W_{t-1}^{q^{\delta}}
\frac{\delta_t(H_t)}{Z_t(H_t)^2}
\{Q_t^{q^{\delta}}(H_t,1)-Q_t^{q^{\delta}}(H_t,0)\}
\{A_t-b_t(H_t)\}.
\end{split}
\end{equation}

On structural-zero/one histories, define the additional term as zero; the unobserved contrast there need not be identified. Assume bounded rewards, fixed finite \(T\), and bounded positive \(\delta\), so the displayed terms are square integrable. At nondegenerate histories the contrasts are identified on support.

\begin{proof}[Derivation of the behavior-adaptive gradient] At a fixed history, \(\partial q_t^{\delta}(1\mid h)/\partial b_t(h)=\delta_t(h)/Z_t(h)^2\), while the derivative for action 0 is its negative. The derivative of policy value with respect to a policy perturbation at \(t\) is its target-occupancy average multiplied by the target-continuation contrast \(Q_t(1)-Q_t(0)\). Relative to observed history occupancy, that average is weighted by \(W_{t-1}^{q^{\delta}}\). Finally, a behavior-score perturbation satisfies \(\dot b_t(h)=E[(A_t-b_t(h))S_{b,t}\mid H_t=h]\). Combining these three derivatives yields the additional covariance term in \eqref{eq:core21}. The fixed-policy derivative is given by Theorem~\ref{core-result-2}. The extra terms are mean-zero action-score components; adding them gives the canonical gradient by the same tangent-space argument. This is the chain-rule construction in Kennedy's general stochastic-intervention lemmas. \end{proof}

A useful check is \(\delta_t\equiv1\). Then \(q=b\) and the adaptive target equals the observed mean return. In \eqref{eq:core21}, the binary action component becomes \(Q_t(H_t,A_t)-V_t(H_t)\); it telescopes with the fixed-policy expression to \(G-E(G)\), the EIF of the observed mean. Omitting the extra term would generally fail this check.

The adaptive target is not conventionally doubly robust against arbitrary misspecification of \(b\), because \(b\) defines the target itself. This document supplies its derivative, not a completed general cross-fitted estimator theorem or implementation for unknown \(b\). The initial empirical program should prioritize the known-randomization case and a finite prespecified \(\delta\) grid. Selecting \(\delta\) continuously on evaluation outcomes needs uniform inference or a separate evaluation sample. No fixed-policy standard error should be silently reused for case 3.

\section{Why static replay and unsupported extrapolation fail}\label{why-static-replay-and-unsupported-extrapolation-fail}

\subsection{A counterexample to holding the future history fixed}\label{a-counterexample-to-holding-the-future-history-fixed}

Let \(T=2\), no baseline covariates, \(L_2=A_1\), and terminal reward \(Y=\mathbf{1}\{A_2=L_2\}\). Consider a recorded trajectory generated by \((A_1,A_2)=(0,0)\), so its recorded intermediate state is \(L_2=0\). The target policy always takes \((1,1)\). A live target execution generates \(L_2=1\) and has reward 1. A static replay that changes the actions to \((1,1)\) but keeps the recorded \(L_2=0\) assigns reward 0. It evaluates a different intervention that freezes a treatment-affected variable.

This example does not require claiming that every replay method is invalid. Restoring a valid pre-decision snapshot and \textbf{executing the new continuation} under the target kernel is a real branch experiment. Also, replay that is correct by a specially justified environment invariance can be valid. What fails is assuming unchanged downstream states without that justification.

Simply copying the recorded outcome is a different operation from recomputing the reward after freezing the state. Outcome copying returns 1 in the example above, just as the live target does, although a policy-independent copied outcome generally cannot distinguish policies. A counterexample must therefore be attached to the specified replay operation and assumptions.

\subsection{Adaptive donor matching can discard the relevant state}\label{adaptive-donor-replay}

Consider another two-decision system, with first action fixed at zero. The first execution reveals \(U\sim\operatorname{Bernoulli}(1/2)\), and visible validation always fails at this decision. The logger then chooses \(A_2\sim\operatorname{Bernoulli}(1/2)\) independently of \(U\). Terminal hidden-test success is \(Y=\ind\{A_2=U\}\); visible validation stops the episode at decision two, including when \(Y=0\). The adaptive target chooses \(A_2=U\), so \(V(\pi)=1\).

Let an unlimited ordered bank contain independent and identically distributed donors from this same task and execution law, with order fixed before observing their records. Define replay to take the first donor's observed state, request the target action there, and return the stored outcome of the earliest donor matching the requested action prefix. Both second actions have a matching donor almost surely; no fallback is needed. If the first donor's logged action equals its state, probability \(1/2\), replay keeps that donor and returns 1. Otherwise replay selects the first later donor whose action equals the \emph{old} state \(U\). Since each donor's action is independent of its own state, selection leaves the replacement state \(U'\) an independent fair bit. Thus
\[
\E(\widehat Y_{\mathrm{replay}})
=\tfrac12\cdot1+\tfrac12\cdot\tfrac12
=\tfrac34,
\qquad V(\pi)=1.
\]
Matching the action prefix has discarded the state that determined the adaptive target's action. The discrepancy occurs with stable kernels, supported randomized actions and unlimited donors at a horizon of two. It does not imply a particular bias on any empirical benchmark.

In the same system, a target with constant second action \(c\) has value \(\P(U=c)=1/2\). Replay selects the first donor with \(A_2=c\), whose state is still a fair bit, and also averages \(1/2\). This narrow positive control explains why replay failure cannot be assumed for every policy. It supplies no general validity theorem for adaptive targets, history-dependent logging probabilities, finite-bank fallback or other forms of state substitution.

\subsection{Nonidentification without support}\label{nonidentification-without-support}

Take one decision and no covariates, with \(A=0\) and \(Y=0\) always in the logs. In world I, the potential reward for action 1 is 0; in world II, it is 1. The observed laws are identical, but the value of ``always action 1'' differs. No estimator based solely on these logs identifies that value. A language-model outcome regressor, a world model, a DR correction, or a different causal vocabulary cannot resolve this observational equivalence without an assumption or new data.

\subsection{What branches can and cannot validate}\label{what-branches-can-and-cannot-validate}

A branch created at a recorded \(H_t\) estimates a continuation effect averaged over the \textbf{branch-selection distribution of histories}, with its own randomized allocation and sampling uncertainty. If histories are selected from behavior trajectories, this is generally not the full target-policy occupancy distribution. Branches can test local calibration of \(Q_t^{\pi}\), reveal downstream divergence, and add experimental support at those particular histories.

To estimate a complete policy value from branches alone, the design must also identify or reweight the target distribution of branch histories and cover histories newly reached under target execution. A finite collection of branches does not repair missing support elsewhere. Fresh root-to-terminal executions under frozen policies provide the clearest external calibration target. Restoring filesystem state, tool sessions, hidden caches, random state, and memory correctly is part of the validity of a branch; text replay alone may not reproduce the same \(H_t\).
